% ch40b-svdq-fp8.tex — SVDQuant(四)：W8A8——per-token E4M3 与低秩残差的融合 GEMM（RFC-S）
% 纯书稿章（不改 cache-dit、不跑 kernel）；源码引用 cache-dit commit fdf42708（PR #1384）。
% 图：全 TikZ inline；listing 走 anchors.yaml 登记的 linerange。
\chapter{SVDQuant（四）：W8A8——per-token E4M3 与低秩残差的融合 GEMM}\label{ch:40b}

\section{本章导读}

第~\ref{ch:38}~章建立了 SVDQuant 的数学主线：通道缩放平滑把激活异常值迁进权重、
低秩分支把平滑后权重里的大奇异方向搬进 BF16、残差只留低比特码；第~\ref{ch:39}~
与~\ref{ch:40}~章把这套机制落到 NVFP4 的 W4A4 实现上。本章回答一个自然的问题：
\textbf{残差码点从 4 bit 换成 8 bit（E4M3），同一套机制还成立吗}？cache-dit 在
PR \#1384（合并提交 \texttt{fdf42708}）里给出的答案是肯定的——SVDQuant-FP8
（W8A8，权重与激活同为 E4M3 加 per-token/per-row 标量尺度）在 SM120 上以两个
融合算子（producer 与 GEMM）落地，且 low-rank 分支的收益结构发生了有趣的变化：rank 128
时 PSNR 相对纯 FP8 提升超过 $5\,\mathrm{dB}$，而 rank 32/64 反而略有回退。

四条主线：

\begin{itemize}[itemsep=1pt,topsep=2pt]
  \item \textbf{扩展的数学条件}（\S\ref{sec:40b-motivation}）：SVDQuant 的分解
        与误差上界对\textbf{任意}反量化映射 $\mathcal Q$ 成立，把 $\mathcal Q$
        从 NVFP4 两级映射换成 per-row E4M3 映射后命题链原样有效；这是"能扩展"
        的全部数学前提，剩下的都是工程。
  \item \textbf{scale 的 K 不变性}（\S\ref{sec:40b-math}）：per-token 激活尺度
        与 per-row 权重尺度都沿归约轴 $K$ 恒定，可以把 scale 完全提出 GEMM
        归约（式~\ref{eq:34-fold-rc} 的特例），因此 W8A8 \textbf{不需要}
        block-scaled mainloop——这是它与 W4A4 在 kernel 结构上更简单的根因；
        但"更简单"不等于"没有新约束"，scale 的施加次序成为一个必须证明的性质。
  \item \textbf{两个融合 kernel}（\S\ref{sec:40b-design}）：producer 在
        一次对激活 tile 的读取里同时产出量化码与 LoRA-down 项 $D$；GEMM 把
        E4M3 主积与 BF16 LoRA-up 修正累进\textbf{同一个} FP32 accumulator，
        全链路没有独立的 BF16 GEMM，也没有对 $[M,N]$ dense 输出的第二次扫描。
  \item \textbf{性能与精度是 trade-off}（\S\ref{sec:40b-perf}）：W8A8 权重
        8 bit、无体积收益，它的定位是\textbf{精度档}：用 $11\%$--$24\%$ 的
        算子级耗时换 rank 128 时 $+5.6\,\mathrm{dB}$ 的 PSNR。
\end{itemize}

\textbf{前置章节}：第~\ref{ch:38}~章（数学主线与记号，必须）、第~\ref{ch:39}~章
（W4A4 的 producer 先例）、第~\ref{ch:40}~章（W4A4 GEMM 与 LoRA 融合先例）、
第~\ref{ch:31}~章（NVFP4 两级 scale 的格式定义）、第~\ref{ch:34}~章（量化 GEMM
的 scale 提出推导）。后续第~\ref{ch:41}~章讲原生 FP8（无低秩分支）的 CUTLASS
Collective 实现，可与本章对照阅读。

\textbf{源码区间与验证方式}：本章讲解 cache-dit 公开仓库 commit
\texttt{fdf42708} 新增的 \texttt{csrc/kernels/svdq/cute/fp8/}（GEMM
\texttt{gemm\_w8a8.*} 与 producer \texttt{quantize\_w8a8.*}）、
\texttt{src/cache\_dit/quantization/svdquant/fp8.py} 与
\texttt{src/cache\_dit/kernels/svdquant\_fp8.py}。纯书稿任务：不改 cache-dit
代码、不运行任何 kernel；数值断言以源码注释与 PR \#1384 的验收记录为准，
\textbf{本书未复测}。完整代码入口见附录~\ref{app:D}。

\section{问题与动机：命题与 $\mathcal Q$ 无关，FP8 残差合法}\label{sec:40b-motivation}

\subsection{从 W4A4 到 W8A8 只差一步映射替换}

回顾第~\ref{ch:38}~章的分解。给定通道缩放 $s\in\mathbb{R}_{>0}^{K}$，
平滑后权重 $\widetilde W = W\cdot\mathrm{diag}(s)$ 被拆成
\begin{equation}
  \widetilde W = L_1 L_2 + R,
  \qquad
  L_1\in\mathbb{R}^{N\times r},\quad L_2\in\mathbb{R}^{r\times K},
  \label{eq:40b-split}
\end{equation}
低秩因子 $L_1$（up）与 $L_2$（down）保留 BF16，残差 $R$ 交给低比特量化器。
第~\ref{ch:38}~章的 Proposition~4.1（转述）给出的误差上界对任意 $X,W$ 与
\textbf{任意}反量化映射 $\mathcal Q$ 成立：W4A4 里 $\mathcal Q$ 是 NVFP4 的
两级映射（E2M1 码字加 16 元素一组的 FP8 micro-scale，第~\ref{ch:31}~章），
W8A8 里它换成
\begin{equation}
  \mathcal Q_{\mathrm{fp8}}(R)_{n,k}
    = \mathrm{wscales}[n]\cdot\bar q_B[n,k],
  \qquad
  \mathrm{wscales}[n]=\frac{\max_k |R[n,:]|}{448},
  \label{eq:40b-wq}
\end{equation}
即每输出通道一个 FP32 尺度的 E4M3 对称量化（空行以 $10^{-12}$ 兜底，
避免零尺度）。映射换了，分解
\eqref{eq:40b-split}、低秩分支对大奇异方向的吸收、
Proposition~4.1 的上界结构\textbf{全部不变}。这就是"能扩展到 FP8"的数学
内容：它不需要新定理，只需要确认既有定理的量词确实覆盖新的 $\mathcal Q$。

工程上真正变化的是两件事，也正是本章的两条技术主线：其一，per-row 权重尺度
与 per-token 激活尺度都沿 $K$ 恒定，scale 可以完全退出 GEMM 归约
（\S\ref{sec:40b-math}）；其二，E4M3 是 8-bit 码点，残差占的字节数是 NVFP4
的约 $1.8$ 倍，W8A8 \textbf{不是压缩方案}——它不省权重显存，定位是花钱买
精度的档位（\S\ref{sec:40b-perf}）。

\subsection{FP8 已经很准，为什么还要低秩分支}\label{sec:40b-why-lora}

这是本章最该想明白的问题。E4M3 的最大有限值是 $448$（$127$ 个正码点），
RN 相对误差上界 $2^{-4}$（第~\ref{ch:27}~章的术语口径：这是误差上界，
binade 下沿的相对 ULP 是 $2^{-3}$），比 E2M1 的 RN 上界 $2^{-2}$ 小
$4$ 倍；看起来 per-token FP8 已经"够用了"。第~\ref{ch:38}~章的通道异常值
分析在这里仍然适用，区别只是打折的程度：激活的异常值固定在少数输入通道上，
per-token 的行 absmax 归一化只是把整行除以同一个数，行内元素之间的
\textbf{相对动态范围}不变，异常通道仍会抬高整行的量化步长、压缩普通通道的
有效级数。位宽从 4 提到 8 只是把"有效级数"的公式
\eqref{eq:38-levels} 里的 $2^b$ 乘了 $16$ 倍（把 $m_j,m$ 取成行内口径
即得行级版本），行内通道悬殊的结构没有变。

更直接的证据来自 PR \#1384 的精度表（\S\ref{sec:40b-perf}，本书未复测）：
rank 32/64 的 PSNR 在 eager 与 compile 四组比较里\textbf{全部}略低于纯 FP8
（$-0.4$--$-1.0\,\mathrm{dB}$），rank 128 才拉开 $+5.2$--$+5.6\,\mathrm{dB}$。
这个排序说明低秩分支要吸收的不是随机码点噪声（那种误差加低秩项只会更差），
而是 per-token FP8 残差里的\textbf{方向性成分}——按 Eckart--Young
（第~\ref{ch:38}~章），吸收方向性残差所需的秩就是残差谱的有效秩。本书的
推测是该有效秩落在百量级（可用 $\mathrm{diag}(\Sigma)$ 的参与比检验），
所以 rank 32/64 只吸走一小部分能量、抵不过码点噪声，rank 128 才跨过门槛；
PR 只给出“rank 128 是精度档”的结论，没有给出这一解释。W8A8 与 W4A4
在低秩分支的用法上由此出现
一个实用分野：W4A4 里 rank 32 是性价比档，W8A8 里 \textbf{rank 128 才是卖点}。

\subsection{W8A8 在 FP8 家族中的位置}

cache-dit 的原生 FP8（第~\ref{ch:41}~章，\texttt{csrc/kernels/fp8/} 族）已经
提供 per-row/per-tensor 粒度的 FP8 GEMM，数学形态是
$\mathrm{out}=(q_Aq_B^{\top})\cdot(\delta_a\delta_b)+\mathrm{bias}$。
W8A8 在它之上只多两样东西：离线阶段的低秩分解 \eqref{eq:40b-split}，与在线
阶段的两条低秩数据通路（$D=X\,(L_2/s)^{\top}$ 与 $D\,L_1^{\top}$）。把它
理解成"带低秩残差修正的 per-token FP8"，比理解成"另一种 FP8 格式"更接近
实现：scale、码点、MMA 全部沿用 FP8 的既有件，新增的只是 producer 里的
LoRA-down 累加与 GEMM 里的 LoRA-up 累加。

\section{数学原理：per-token 量化、scale 的 $K$ 不变性与施加次序}\label{sec:40b-math}

\subsection{记号与量化三式}

沿用第~\ref{ch:38}~章的符号方向：$X\in\mathbb{R}^{M\times K}$、
$W\in\mathbb{R}^{N\times K}$、$Y=XW^{\top}$；$L_1\in\mathbb{R}^{N\times r}$
是 up 因子、$L_2\in\mathbb{R}^{r\times K}$ 是 down 因子
（注意与论文原文的转置约定，第~\ref{ch:38}~章已交代）。实现里
\texttt{lora\_down} 存的是 $L_2/s$（已折入逆平滑）、\texttt{lora\_up} 存的
是 $L_1$，与本书记号一致。

\textbf{离线权重侧}。分解 \eqref{eq:40b-split} 之后，残差按
\eqref{eq:40b-wq} 量化；低秩因子里 $L_2$ 要除以平滑向量折回原始量纲：
\begin{equation}
  \texttt{lora\_down} = L_2 / s,
  \qquad
  \texttt{lora\_up} = L_1,
  \label{eq:40b-offline}
\end{equation}
这个除法在导出时用 float64 完成（与 W4A4 导出器同一约定），漏掉它整个低秩
分支会差一个因子 $s$。

\textbf{在线 producer 侧}。对每一行激活 $X_{m,:}$，先做逐输入通道平滑
$X/s$，再取行 absmax 定 per-token 尺度，量化码用 satfinite 舍入（溢出钉在
E4M3 最大有限值而不是刷成 inf）：
\begin{equation}
  \mathrm{sa}[m]
    = \frac{\max\bigl(\mathrm{amax}_k |X[m,k]/s[k]|,\ \varepsilon\bigr)}{448},
  \qquad
  \bar q_A[m,k]
    = \mathrm{satfinite\_e4m3}\!\left(\frac{X[m,k]/s[k]}{\mathrm{sa}[m]}\right),
  \label{eq:40b-aq}
\end{equation}
$\varepsilon=10^{-12}$ 防止全零行产生零尺度。同一趟对激活 tile 的读取还要
产出低秩 down 项：
\begin{equation}
  D[m,r] = \mathrm{BF16}\Bigl(\sum_{k} X[m,k]\cdot\texttt{lora\_down}[r,k]\Bigr),
  \label{eq:40b-d}
\end{equation}
\textbf{$D$ 消耗的是原始激活 $X$ 而不是平滑后的 $X/s$}。因为
\texttt{lora\_down} 里已经带着 $1/s$，乘积
$X\,(L_2/s)^{\top}=(X/s)\,L_2^{\top}=X\,\mathrm{diag}(1/s)\,L_2^{\top}$
直接就是低秩分支在原始量纲下的结果。这是本章
最容易写错的一处数据流：量化链与低秩链\textbf{消费的不是同一份数值}，只是
共用同一次访存。

\textbf{GEMM 侧}。两个乘积共用同一个 FP32 累加器，scale 与 bias 在它们之间
施加，最后只舍入一次：
\begin{equation}
  \mathrm{out}[m,n]
    = \mathrm{BF16}\Bigl(
        \underbrace{\Bigl(\sum_k \bar q_A[m,k]\,\bar q_B[n,k]\Bigr)}_{P_{m,n}}
        \cdot \mathrm{sa}[m]\,\mathrm{wscales}[n]
        + \mathrm{bias}[n]
        + \underbrace{\sum_r D[m,r]\,L_1[n,r]}_{G_{m,n}}
      \Bigr).
  \label{eq:40b-gemm}
\end{equation}

\subsection{scale 可以提出，但次序不能乱}

式~\eqref{eq:40b-gemm} 里 $\mathrm{sa}[m]$ 只随行、$\mathrm{wscales}[n]$ 只
随列，两者都沿 $K$ 恒定，于是
\begin{equation}
  \sum_k (\bar q_A\cdot\mathrm{sa})(\bar q_B\cdot\mathrm{wscales})
    = \mathrm{sa}\cdot\mathrm{wscales}\sum_k \bar q_A\,\bar q_B,
  \label{eq:40b-fold}
\end{equation}
这正是第~\ref{ch:34}~章式~\ref{eq:34-fold-rc} 的"scale 完全出括号"与
第~\ref{ch:41}~章式~\ref{eq:41-fp8-evt} 的 EVT 施加在 W8A8 语境下的重述：
$\mathrm{sa}$/$\mathrm{wscales}$ 就是那两个式子里的 $s_A$/$s_B$（亦即
第~\ref{ch:34}~章的 $\delta_{a,m}$/$\delta_{b,n}$）。
由此得到与 W4A4 的结构对照：NVFP4 的 micro-scale 是 16 元素一组、\textbf{沿
$K$ 变化}，提不出归约，只能在 mainloop 里逐 block 乘（block-scaled MMA，
第~\ref{ch:40}~章）；W8A8 的 scale 是行/列级、\textbf{沿 $K$ 不变}，整个
GEMM mainloop 是普通的 FP8 MMA，scale 全部推迟到 epilogue。这是 FP8 扩展在
kernel 上更简单的根因，也是图~\ref{fig:40b-format} 想传达的本质：
\textbf{差别的本质不是 scale 的层数，而是 scale 是否沿归约轴变化}。

出括号之后还剩一个 W4A4 没有的问题：式~\eqref{eq:40b-gemm} 里 scale 乘在
主积 $P$ 上，却\textbf{不能}乘到低秩项 $G$ 上。验证只差一行代数：若把次序
换一下、先累加再统一乘尺度，得到的会是
\begin{equation}
  \mathrm{sa}\cdot\mathrm{wscales}\cdot(P+G)
    \;=\; \underbrace{\mathrm{sa}\,\mathrm{wscales}\,P}_{\text{想要的}}
      + \underbrace{\mathrm{sa}\,\mathrm{wscales}\,G}_{\text{多出来的}},
  \label{eq:40b-leak}
\end{equation}
误差项是 $(\mathrm{sa}\cdot\mathrm{wscales}-1)\,G$。$\mathrm{sa}$ 与
$\mathrm{wscales}$ 都是 absmax 除以 448 的量级：残差 amax 远小于 448 是
常态（低秩分支刚把大奇异方向搬走），于是 $\mathrm{sa}\cdot\mathrm{wscales}$
通常\textbf{远小于} 1，后乘会把低秩项整体压到接近零，效果上等于丢掉 $G$、
退化为纯 FP8；只有两侧 amax 都接近 448 的极端情形才可能大于 1 而放大。
W4A4 没有这个问题是因为它的低秩项通过 LoRA-first 的代数补偿在
\textbf{同一个尺度约定}里归一（第~\ref{ch:40}~章式~\ref{eq:40-lorafirst}）；
在本 kernel 的单累加器设计下，W8A8 的低秩项直接以原始量纲参与，次序因此
固定为先乘 scale 再累加；等价的绕法是先把 $G$ 预除 $1/(\mathrm{sa}\cdot
\mathrm{wscales})$、最后再乘回，代价是每元素一次除法。实现里这条约束由
函数调用顺序固化：
\texttt{apply\_scales}（源码 gemm\_w8a8.cuh L261--286）就地改写 FP32
fragment，\texttt{gemm\_bf16\_lora}（L343--378）随后把 LoRA-up 积累进
\textbf{同一个} fragment：两个 MMA 的 C partition 在编译期被
static\_assert 锁成逐元素一致（L97--102），所以"累进同一累加器"在布局层面
是定理而不是巧合。

\subsection{与 W4A4 的三点区别}\label{sec:40b-compare}

\begin{table}[H]
\centering\small
\caption{SVDQuant 的两种残差格式对照。W4A4 列为第~\ref{ch:39}--\ref{ch:40}~章的
现状；W8A8 列为本章（cache-dit \texttt{fdf42708}）。}
\label{tab:40b-compare}
\begin{tabular}{@{}p{2.6cm}p{5.2cm}p{5.2cm}@{}}
\toprule
 & W4A4（NVFP4） & W8A8（E4M3） \\
\midrule
残差码点 & E2M1，4 bit，两码点 packed 一字节 & E4M3，8 bit，一码点一字节 \\
scale 结构 & UE4M3 micro（16 元素一组）+ FP32 标量 $\alpha$，
             视 runtime 另加逐列 \texttt{wcscales} & 每行（激活）/每输出通道
             （权重）一个 FP32 尺度，共一层 \\
\textbf{scale 是否沿 $K$ 变化} & \textbf{是}（micro-scale 按组沿 $K$ 分布） &
             \textbf{否}（行/列级标量） \\
scale 施加位置 & mainloop 内逐 block（block-scaled MMA） &
             epilogue，就地在 FP32 fragment 上 \texttt{fmaf} \\
激活量化 & NVFP4 per-token 两级（第~\ref{ch:39}~章） & E4M3 per-token 单级
             （式~\eqref{eq:40b-aq}） \\
权重视显存 & 约 $4.5$ bit/元素（码点 4 bit 加 scale 分摊） & $8$ bit/元素
             （加 scale 可忽略） \\
低秩分支收益档 & rank 32 即性价比档 & rank 128 才是精度档 \\
checkpoint 序列化 & 支持（native pack 格式） & 显式拒绝（格式未定型） \\
\bottomrule
\end{tabular}
\end{table}

三点区别按重要性排序。第一，\textbf{scale 与 $K$ 的关系}：这是 kernel 结构
差异的总开关——沿 $K$ 不变的 scale 可以退出 mainloop，沿 $K$ 变化的必须
留在 mainloop。第二，\textbf{量化粒度与格式复杂度}：W4A4 要维护两级 scale
的 pack 布局、SFA/SFB 的 smem atom 与 block-scaled MMA 指令（第~\ref{ch:40}~章
的大半篇幅）；W8A8 的 producer 与 GEMM  layouts 都是教科书式的行主序
E4M3，尺度是两个长度 $M$/$N$ 的 FP32 向量。第三，\textbf{经济模型}：W4A4 的
卖点是 $4.5$ bit/元素的权重视显存，W8A8 权重视显存约为其 $1.8$ 倍、没有
体积收益，它存在的理由是精度表上 rank 128 的那一列（\S\ref{sec:40b-perf}）。
换言之，W4A4 回答"怎么把 Linear 变便宜"，W8A8 回答"per-token FP8 还不够
准时怎么修"。

图~\ref{fig:40b-format} 把两个格式的位布局画在一起。读图注意：NVFP4 侧
画的是 micro 层（16 元素一组的 E4M3 尺度）与 global 层两层，W8A8 侧只有
per-row 一层；但本节的论点不依赖"几层"，而依赖"沿 $K$ 是否变化"；即便
把 W4A4 的 global 层去掉只剩 micro 层，它的 scale 依旧沿 $K$ 变化、依旧
需要 block-scaled mainloop。

\section{设计与实现}\label{sec:40b-design}

\subsection{fused producer：一次读取，两种产物}\label{sec:40b-producer}

producer 的数学契约是式~\eqref{eq:40b-aq} 与~\eqref{eq:40b-d}：从 BF16
激活 $X$ 产出 E4M3 码 $\bar q_A$、FP32 per-token 尺度 sa、BF16 低秩项 $D$。
cache-dit 把它拆成两个 kernel 共享一个入口（quantize\_w8a8.cuh），拆分的
理由是并行度：

\begin{itemize}[itemsep=1pt,topsep=2pt]
  \item \textbf{amax kernel}（L96--125）：\textbf{每行一个 CTA}，CTA 内 8 个
        warp 切 $K$，每个线程以 16 B（8 个 BF16）为单位向量化读取，先做
        逐通道平滑 $X/s$ 再取绝对值，warp 内 \texttt{shfl\_xor} 归约出行
        amax，最后 $\mathrm{sa}=\max(\mathrm{amax},\varepsilon)/448$ 写入
        gmem。行级 grid 提供足够多的独立 load 流来掩 HBM 延迟，
        128 行一 tile 的 CTA 自己做不到这一点（源码注释 L93--95 明说这是在
        镜像 cache-dit 原生 FP8 producer 的并行分解）。
  \item \textbf{主 kernel}（L127--292）：CTA 拥有 128 行与 $K$ 的一个连续
        分片（grid.y 是 split 维，分片大小由
        \texttt{w8a8\_\allowbreak{}quant\_\allowbreak{}k\_\allowbreak{}splits}
        决定，目标让总 CTA 数 $\ge256$），
        内部按 64 宽的 K tile 流水：先把当前 tile 从寄存器写进 smem，
        随即发起 tile $kt{+}1$ 的寄存器预取（\texttt{x\_regs}/\texttt{l\_regs}，
        各两份），屏障之后 LDSM 读出做 BF16 MMA 累加 $D$、再在
        \textbf{同一份寄存器副本}上逐元素 $X/s$、乘预计算的
        $\mathrm{sa}^{-1}$、satfinite 转成 E4M3 写回 gmem——预取与
        MMA、量化两步重叠。
\end{itemize}

amax kernel 本体（L96--125）：向量化的 8 元素读取、平滑后取绝对值、warp
内 \texttt{shfl\_xor} 归约、最后一行除以 448 得 sa。

\begin{lstinputlisting}[linerange={96-125},firstnumber=auto,
  title={quantize\_w8a8.cuh L96--125（amax kernel：每行一 CTA 的 per-token 尺度）}]{../../../../cache-dit/csrc/kernels/svdq/cute/fp8/quantize/quantize_w8a8.cuh}
\end{lstinputlisting}

量化与 $D$ 的累加共用同一批寄存器副本，是这段代码的关键安排：$X$ 从
gmem 只读一次（16 B 向量化），smem 只作 MMA 的暂存，E4M3 写回按
8 码点 8 B 一次写出、与读取同构地 coalesced。per-token 尺度用
\textbf{预计算倒数}（CTA 开头把 128 行的 $1/\mathrm{sa}$ 放进 smem，
L154--155），量化循环里对 sa 的除法因此变成乘法；逐输入通道的平滑
（$X/s[k]$）仍是一次除法。源码注释（L59--60）特意交代：倒数乘法与
fast-math 除法可能让 E4M3 翻一个 ULP，所以转换指令固定用 RTN 加 satfinite 的
\texttt{\_\_nv\_cvt\_float2\_to\_fp8x2}，与 \texttt{csrc/kernels/fp8} 的
producer 逐指令一致。

$D$ 的累加是 SM80 的 \texttt{16x8x16} BF16 MMA（E4M3 化之前先算出低秩项，
两者都以 BF16/FP32 语义工作），全 $K$ 循环只在 FP32 累加器里转，最后
\textbf{单次}舍入到 BF16；若 $K$ 被切成多个 split，各 split 的 FP32
partial fragment 用 relaxed 语义的 \texttt{atomicAdd} 归并进 scratch，
由一个 finalize kernel 做那次唯一的舍入（L275--303）。这套"并行度换
可复现性"的取舍与第~\ref{ch:40}~章 split-K 的约定一致：原子加不保证
逐位可复现，$D$ 的等价性判断用误差门槛而不是 bitwise。scratch 不在
Python 签名里——pybind 层（svdq.cpp 的 \texttt{quantize\_w8a8}，见
\texttt{fdf42708} 的 L214--235）检测到 $\mathrm{rank}>0$ 且 splits $>1$ 时用
\texttt{at::zeros} 分配 $[M,\mathrm{rank}]$ 的 FP32 缓冲，Python 侧完全不可见。
（附录~\ref{app:C} 的混合 checkout 里 \texttt{svdq.cpp} 取自
\texttt{eced0158}，不含该函数；查这段请直接看 \texttt{fdf42708}。）

两段关键代码。先是主 kernel 的处理核（process\_tile 上半，L213--246）：
寄存器到 smem 的 staging、下一片的预取、LDSM 与 BF16 MMA 都发生在这里。

\begin{lstinputlisting}[linerange={213-246},firstnumber=auto,
  title={quantize\_w8a8.cuh L213--246（process\_tile 上段：ST smem、双缓冲预取与 $D$ 的 MMA）}]{../../../../cache-dit/csrc/kernels/svdq/cute/fp8/quantize/quantize_w8a8.cuh}
\end{lstinputlisting}

再是量化写回（process\_tile 下半，L247--266）：同一份
\texttt{x\_regs} 副本上完成 $X/s$、乘倒数、satfinite 转换。

\begin{lstinputlisting}[linerange={247-266},firstnumber=auto,
  title={quantize\_w8a8.cuh L247--266（process\_tile 下段：量化写回）}]{../../../../cache-dit/csrc/kernels/svdq/cute/fp8/quantize/quantize_w8a8.cuh}
\end{lstinputlisting}

split-K 的写回与调度函数（L275--312）：

\begin{lstinputlisting}[linerange={275-312},firstnumber=auto,
  title={quantize\_w8a8.cuh L275--312（split-K 原子写回、finalize 与 k\_splits 目标）}]{../../../../cache-dit/csrc/kernels/svdq/cute/fp8/quantize/quantize_w8a8.cuh}
\end{lstinputlisting}

注意 \texttt{w8a8\_quant\_k\_splits} 的目标公式（L307--312）：以
$256$ 个 CTA 为目标反推 split 数，且不超过 K tile 总数；rank 0 没有
$D$ 要归约，总是按该目标选 split（L332--334 在 dispatcher 处决定；
K 很小时 split 退化为 1）。host 入口
的校验与 rank 分派（quantize\_w8a8.cu）：$M$ 是 128 的倍数、$K$ 是 64 的
倍数、rank $\in\{0,32,64,128\}$，rank 32/64/128 各一个实例化 TU 以并行
编译。

\begin{lstinputlisting}[linerange={52-64},firstnumber=auto,
  title={quantize\_w8a8.cu L52--64（host 几何与 dtype 约束）}]{../../../../cache-dit/csrc/kernels/svdq/cute/fp8/quantize/quantize_w8a8.cu}
\end{lstinputlisting}

\subsection{W8A8 GEMM：双乘积单累加器}\label{sec:40b-gemm}

GEMM 的数学契约是式~\eqref{eq:40b-gemm}。结构上它把第~\ref{ch:40}~章
W4A4 GEMM 的两条原则原样继承下来——\textbf{没有任何独立 BF16 GEMM、
没有对 dense 输出的二次扫描、FP32$\to$BF16 全 kernel 只发生一次}——
但 mainloop 本身反而\textbf{更常规}：没有 block-scaled MMA、没有 SFA/SFB
布局、没有 G2R 逆映射，只有 TMA 加 LDSM 加 SM120 FP8 MMA 的标准流水。
新增的复杂度全部集中在 LoRA 分支与它的调度上。

\textbf{traits 与双 MMA 布局契约}（gemm\_w8a8.cuh L84--102）。FP8 主积用
SM120 的 \texttt{16x8x32} E4M3 MMA atom，AtomLayout $4\times2\times1$、
Permutation $128\times32\times32$；LoRA-up 用 SM80 的
\texttt{16x8x16} BF16 MMA atom，M/N 完全相同、只有 K 维换成
\texttt{kRankTile}$=32$。两个 tiled MMA 的线程数与 C partition 被两条
static\_assert 锁死成逐类型相等：这就是两个乘积能共享
\texttt{tCrC} 的布局前提，第~\ref{ch:40}~章在 W4A4 里建立过同一结论。

\begin{lstinputlisting}[linerange={84-102},firstnumber=auto,
  title={gemm\_w8a8.cuh L84--102（E4M3/BF16 双 MMA atom 与共享累加器的 static\_assert）}]{../../../../cache-dit/csrc/kernels/svdq/cute/fp8/gemm/gemm_w8a8.cuh}
\end{lstinputlisting}

\textbf{smem 预算}（L127--144）。主流水两 stage，每 stage A（$128\times128$
E4M3）16 KiB 加 B 16 KiB；LoRA 专用 ring 也是两 stage，每 stage D
（$128\times32$ BF16）8 KiB 加 $L_1$ 8 KiB；合计 96 KiB，被
\texttt{kSmemBytes} $\le 99$ KiB 的静态断言守住。输出 staging tile
（$128\times128$ BF16，32 KiB）不复用 LoRA ring，而是叠在主流水的存储上
（断言 L138--139 要求 \texttt{SmemLayoutO} 字节数不超过主流水总量）；
epilogue 跑时主流水已排空，这是第~\ref{ch:40}~章 overlay 手法的简化版。
\textbf{注意实例化参数}：cuh 的模板形参默认 \texttt{kStages\_=3}，但 host
只实例化 \texttt{launch\_gemm\_w8a8<128,2,128>}（gemm\_w8a8.cu L76）；
smem 的账只按两 stage 成立，读代码以实例化为准。

\begin{lstinputlisting}[linerange={127-144},firstnumber=auto,
  title={gemm\_w8a8.cuh L127--144（两 stage 主流水加 LoRA ring 的 96 KiB 预算与 overlay 断言）}]{../../../../cache-dit/csrc/kernels/svdq/cute/fp8/gemm/gemm_w8a8.cuh}
\end{lstinputlisting}

\textbf{主循环}（L226--246）没有任何 scale：TMA barrier 等到 stage 满，
LDSM 逐 32 宽 k-block 读 fragment，MMA 写 \texttt{tCrC}，warp leader
释放 stage。

\begin{lstinputlisting}[linerange={226-246},firstnumber=auto,
  title={gemm\_w8a8.cuh L226--246（gemm\_fp8\_main：无 scale 的 E4M3 主循环）}]{../../../../cache-dit/csrc/kernels/svdq/cute/fp8/gemm/gemm_w8a8.cuh}
\end{lstinputlisting}

\textbf{scale 施加}（L261--286）是 §\ref{sec:40b-math} 那条次序约束的代码
形态：\texttt{acc = fmaf(acc * sa[m], sw[n], bias[n])} 就地改写
fragment，发生在 LoRA 累加\textbf{之前}。scale 的寻址依赖一个 regroup
助手 \texttt{w8a8\_acc\_rowcol}（L32--40）：SM120 atom 的 C 布局把两行
交错进一对寄存器，先把 fragment 重排成 (row, col) 两个独立 mode，才能
按行取 sa、按列取 sw 与 bias。

\begin{lstinputlisting}[linerange={261-286},firstnumber=auto,
  title={gemm\_w8a8.cuh L261--286（apply\_scales：先乘 sa 再 LoRA 累加的次序固化）}]{../../../../cache-dit/csrc/kernels/svdq/cute/fp8/gemm/gemm_w8a8.cuh}
\end{lstinputlisting}

\textbf{LoRA-up 累加}（L343--378）按 32 宽的 rank 切片循环（与 FP4 kernel
同粒度，每切片 $D$ 8 KiB、$L_1$ 8 KiB）。调度上做一个低开销重叠：
consumer 的第一个 warp（相对编号 \texttt{tid<32}）在 mainloop 开始前
用 \texttt{cp.async} 预取前 $\min(\mathrm{nRt},2)$ 个切片填 LoRA ring，
与主循环并行；rank 128 时 $\mathrm{nRt}=4$，超出预取 ring 的下标 2、3
（第 3、4 片）
在循环内用 NamedBarrier 保护按需 refill（L363--373）。SW64 的 smem atom
让 LDSM 读 32 宽 BF16 切片零 bank conflict（源码注释 L111--113：每行
64 B 恰是一个 atom）。

\begin{lstinputlisting}[linerange={343-378},firstnumber=auto,
  title={gemm\_w8a8.cuh L343--378（gemm\_bf16\_lora：rank 切片累加与 rank 128 的按需 refill）}]{../../../../cache-dit/csrc/kernels/svdq/cute/fp8/gemm/gemm_w8a8.cuh}
\end{lstinputlisting}

\textbf{warp specialization 与 epilogue}。kernel 共 384 线程：
\texttt{tid}<128 是 producer（\texttt{warpgroup\_reg\_dealloc<32>} 让出
寄存器，只有 tid 0 发 TMA，发完即 return）；其余 256 线程是 consumer
（\texttt{warpgroup\_reg\_alloc<224>}）。consumer 的调用链一句话：
预取 LoRA 切片 $\to$ \texttt{gemm\_fp8\_main} $\to$ \texttt{apply\_scales}
$\to$ \texttt{gemm\_bf16\_lora} $\to$ 单次 FP32$\to$BF16 $\to$ STSM 进
smem $\to$ TMA store。

\begin{lstinputlisting}[linerange={507-518},firstnumber=auto,
  title={gemm\_w8a8.cuh L507--518（producer 分派：reg\_dealloc 与 TMA 发射环）}]{../../../../cache-dit/csrc/kernels/svdq/cute/fp8/gemm/gemm_w8a8.cuh}
\end{lstinputlisting}

\begin{lstinputlisting}[linerange={520-545},firstnumber=auto,
  title={gemm\_w8a8.cuh L520--545（consumer 调用链：LoRA 预取到 TMA store 的五步）}]{../../../../cache-dit/csrc/kernels/svdq/cute/fp8/gemm/gemm_w8a8.cuh}
\end{lstinputlisting}

单次舍入的证据在 \texttt{stage\_output\_tile}（L388--405）：
\texttt{NumericArrayConverter} 对整个 fragment 只做一次
FP32$\to$BF16，之后的 STSM 与 TMA store 都搬的是 BF16。

\begin{lstinputlisting}[linerange={388-405},firstnumber=auto,
  title={gemm\_w8a8.cuh L388--405（stage\_output\_tile：全 kernel 唯一一次 FP32 到 BF16 转换）}]{../../../../cache-dit/csrc/kernels/svdq/cute/fp8/gemm/gemm_w8a8.cuh}
\end{lstinputlisting}

host 侧（gemm\_w8a8.cu）先做完整校验（E4M3 码、BF16 输出与 LoRA 张量、
FP32 尺度、$M/N/K$ 皆 128 的倍数、rank $\in\{0,32,64,128\}$），再取
实例化 \texttt{launch\_gemm\_w8a8<128,2,128>}；TMA 描述符按
"指针 $+$ shape $+$ rank" 作 key 缓存（上限 32 条，超出清空），因为描述符
只是拷贝配方。

\begin{lstinputlisting}[linerange={59-83},firstnumber=auto,
  title={gemm\_w8a8.cu L59--83（host 几何约束与 $<$128,2,128$>$ 实例化）}]{../../../../cache-dit/csrc/kernels/svdq/cute/fp8/gemm/gemm_w8a8.cu}
\end{lstinputlisting}

\subsection{Python 量化流程与显式拒绝}\label{sec:40b-python}

算子层之上，cache-dit 提供与 W4A4 平行的打包模块
\texttt{SVDQW8A8Linear}（文件 \texttt{fp8.py}，在
\texttt{src/cache\_dit/} 下的 \texttt{quantization/svdquant/} 中）。
buffers 按 kernel 的 128 宽几何 pad（输入与输出宽度各向上取 128 的
倍数，即 \texttt{quantized\_in\_features} 与
\texttt{quantized\_out\_features} 两个量），保留原始宽度供 forward
切片回去；\texttt{dequantize\_weight} 按 $(R+L_1L_2)/s$（即
$R/s+L_1(L_2/s)$）重建权重，是核对残差、尺度与低秩因子三者约定
没有互相拧的现成钩子。

\begin{lstinputlisting}[linerange={139-147},firstnumber=auto,
  title={fp8.py L139--147（残差 per-row 量化与 proj\_down 折逆平滑）}]{../../../../cache-dit/src/cache_dit/quantization/svdquant/fp8.py}
\end{lstinputlisting}

forward 走完整融合链：pad 到 128 行 $\to$ \texttt{svdq\_quantize\_w8a8}
（producer，\texttt{pad\_rows=True} 时保留 pad 供 GEMM 直接消费）
$\to$ \texttt{svdq\_gemm\_w8a8} $\to$ 切回 \texttt{out\_features}。
两个算子在 \texttt{svdquant\_fp8.py} 里注册成
\texttt{cache\_dit\_cuda\_ops} 下的 torch.library 算子，CUDA impl 与
fake 成对出现，fake 的 shape/dtype 与 impl 完全一致，这是
torch.compile 全图追踪不被打断的前提。

\begin{lstinputlisting}[linerange={48-71},firstnumber=auto,
  title={svdquant\_fp8.py L48--71（GEMM 算子的 define / CUDA impl / register\_fake）}]{../../../../cache-dit/src/cache_dit/kernels/svdquant_fp8.py}
\end{lstinputlisting}

\begin{lstinputlisting}[linerange={178-204},firstnumber=auto,
  title={fp8.py L178--204（SVDQW8A8Linear.forward：pad、融合 producer+GEMM、切片返回）}]{../../../../cache-dit/src/cache_dit/quantization/svdquant/fp8.py}
\end{lstinputlisting}

PTQ 入口 \texttt{quantize\_svdq\_fp8} 与 W4A4 共用校准框架：非 DQ 路径
要求 \texttt{calibrate\_fn} 并通过 forward pre-hook 捕获逐层激活 span，
\texttt{smooth\_strategy} 四分支（\texttt{activation} 用校准数据算
SmoothQuant 因子，\texttt{identity} 不平滑，\texttt{weight}/\texttt{
weight\_inv} 是无数据的 DQ 导出）。

\textbf{显式拒绝清单}（比"支持什么"更能说明设计边界）：
\begin{itemize}[itemsep=1pt,topsep=2pt]
  \item checkpoint 序列化：\texttt{serialize\_to} 非 None 即
        \texttt{NotImplementedError}：W4A4 的 save/load 格式描述的是
        NVFP4 pack 布局，与 E4M3 残差不兼容，宁可拒绝也不静默写错；
  \item \texttt{few\_shot} 平滑策略：显式 \texttt{ValueError}；
  \item SVDQ 通用校验（config.py）：\texttt{components\_to\_quantize}、
        \texttt{precision\_plan}、\texttt{per\_tensor\_fallback} 三项不支持；
        FP8 精度分支另要求 \texttt{svdq\_kwargs} 里\textbf{保持默认值}的恰有
        11 个键：\texttt{runtime\_kernel}、\texttt{lora\_layout}、
        \texttt{fused\_mlp}、\texttt{layerwise\_offload}、
        \path{offload_quantized_layers_to_cpu}、
        \path{defer_move_to_execution_device}、
        \path{few_shot_auto_compile}、
        \texttt{svd\_refine\_iters}、\texttt{streaming}、
        \path{activation_buffer_flush_sample_count}、
        \path{activation_buffer_flush_cpu_bytes}，
        任一取非默认值即 \texttt{ValueError}；其余键
        （\texttt{calibrate\_precision}、\texttt{quantize\_device}、
        \texttt{smooth\_strategy} 等）仍按通用校验生效；
  \item 非 BF16 的 \texttt{torch\_dtype}：模块构造即拒绝。
\end{itemize}
PR \#1384 的说明文字还把 tensor parallel 列入"不支持"，但代码层没有对应
校验（\texttt{svdquant/dtensor.py} 提供的是 W4A4 的 TP 包装）；按"每条
断言有代码锚点"的纪律，本章不把它写成代码断言，只记为 PR 文档的口径。

\section{图示}

本章设计图按四个问题读：数据\textbf{从哪来}（图~\ref{fig:40b-pipeline}：
一次读取双产物）、\textbf{格式差在哪}（图~\ref{fig:40b-format}：scale
是否沿 $K$ 变化）、\textbf{什么时候施加 scale}（图~\ref{fig:40b-acc-fusion}
的 accumulator 时序）、\textbf{代价是什么}（图~\ref{fig:40b-tradeoff}的
精度--耗时平面）。

\begin{figure}[H]
\centering
\begin{tikzpicture}[font=\scriptsize,
  box/.style={draw, rounded corners=2pt, align=center, inner sep=4pt},
  wide/.style={box, text width=5.0cm, minimum height=3.6cm},
  narr/.style={box, text width=1.7cm, minimum height=1.1cm},
  arr/.style={->, thick}]
  \node[narr, fill=tkGreen!18] (x) at (0.55,2.4) {$X$ $[M,K]$\\BF16};
  \node[wide, fill=tkOrange!18] (prod) at (4.7,2.4)
    {\textbf{fused producer（两个 launch）}\\[3pt]
     \texttt{amax} kernel：每行一 CTA，行 absmax $\to$ sa\\
     主 kernel：CTA $=$ 128 行 $\times$ $K$ 分片（split-K）\\[2pt]
     smem 写入 $\to$ 预取 tile $kt{+}1$ $\to$ LDSM\\
     \quad$\circ$ BF16 MMA 累加 $D$（FP32，单次舍入）\\
     \quad$\circ$ 同一份 \texttt{x\_regs}：$X/s\to\times\,\mathrm{sa}^{-1}$\\
     \qquad$\to$ satfinite E4M3 写回（8 码点 8 B）};
  \node[wide, fill=tkBlue!18] (gemm) at (11.2,2.4)
    {\textbf{W8A8 GEMM（一个 launch，384 线程）}\\[3pt]
     producer 128T：TMA 发射 A/B（tid 0，reg\_dealloc 32）\\
     consumer 256T（reg\_alloc 224）：\\[2pt]
     \quad$\circ$ E4M3 mainloop：TMA$\to$LDSM$\to$MMA $\to$ $P$\\
     \quad$\circ$ \texttt{apply\_scales}：$P\!\to\!\mathrm{sa}\!\cdot\!\mathrm{sw}P{+}\mathrm{bias}$\\
     \quad$\circ$ LoRA 切片（32 宽）BF16 MMA 累加 $G$\\
     \quad$\circ$ 单次 FP32$\to$BF16 $\to$ STSM $\to$ TMA store};
  \node[narr, fill=tkRed!18] (out) at (15.05,2.4) {out $[M,N]$\\BF16};
  \draw[arr] (x) -- (prod);
  \draw[arr] (prod) -- (gemm)
    node[midway, above=2pt, font=\tiny] {$\bar q_A$, sa, $D$};
  \draw[arr] (gemm) -- (out);
  \node[anchor=north west, align=left, text width=15.6cm, font=\tiny] at (-0.45,-0.6)
    {贯穿全链的三条不变量：量化链与低秩链消费\textbf{不同的数值}（$X/s$ 与 $X$，
     靠 lora\_down 里的 $1/s$ 对齐量纲）但共用同一次访存；$D$ 与 $\bar q_A$ 同kernel产出，
     \textbf{没有独立 BF16 GEMM}；LoRA 修正累进同一 FP32 fragment、输出只舍入一次，
     \textbf{没有对 $[M,N]$ dense 输出的二次扫描}。};
\end{tikzpicture}
\caption{W8A8 端到端数据流。producer 一次读取产出三种产物；GEMM 把 E4M3 主积、
scale、bias 与低秩修正排进同一个 FP32 累加器的生命周期的四段。}
\label{fig:40b-pipeline}
\end{figure}

\begin{figure}[H]
\centering
\begin{tikzpicture}[font=\scriptsize,
  byte/.style={draw, minimum height=0.62cm, inner sep=1.5pt, align=center},
  brace/.style={decorate,decoration={brace,amplitude=4pt}}]
  % NVFP4 行
  \node[byte, fill=tkOrange!18, minimum width=9.2cm, anchor=west] (codes)
    at (0.2,2.6) {E2M1 码点流：两码点 packed 一字节，$4$ bit/元素};
  \node[byte, fill=tkOrange!35, minimum width=3.1cm, anchor=west] (micro)
    at (9.9,2.6) {UE4M3 micro-scale，每 16 元素一个（沿 $K$ 分布）};
  \node[byte, fill=tkOrange!50, minimum width=2.2cm, anchor=west] (glob)
    at (13.6,2.6) {FP32 global/逐列层};
  \node[anchor=west, font=\scriptsize\bfseries] at (0.2,3.35) {W4A4（NVFP4，两级至三级）};
  \draw[tkOrange!75!black, thick]
    (9.9,2.24) -- (15.8,2.24) node[midway, below=3pt, font=\tiny, text=tkOrange!75!black]
    {scale 沿 $K$ 变化 $\Rightarrow$ 提不出归约，留在 block-scaled mainloop};
  % W8A8 行
  \node[byte, fill=tkBlue!18, minimum width=9.2cm, anchor=west] (acodes)
    at (0.2,0.9) {E4M3 码点流：一码点一字节，$8$ bit/元素};
  \node[byte, fill=tkBlue!35, minimum width=6.2cm, anchor=west] (as)
    at (9.95,0.9) {FP32 sa$[m]$（每 token）与 sw$[n]$（每输出通道）};
  \node[anchor=west, font=\scriptsize\bfseries] at (0.2,1.65) {W8A8（E4M3，一层）};
  \draw[tkBlue!70!black, thick]
    (9.95,0.54) -- (16.15,0.54) node[midway, below=3pt, font=\tiny, text=tkBlue!70!black]
    {scale 沿 $K$ 不变 $\Rightarrow$ 提出归约，epilogue 一次 fmaf};
\end{tikzpicture}
\caption{两种残差格式的位布局对照。W4A4 侧画出 micro 层与 global/逐列层
（两级，另加 runtime 可选的逐列层）
（W4A4 的 runtime 差异见第~\ref{ch:40}~章）；W8A8 侧只有行/列级一层。
本图只对照 micro 层与 per-row 层的粒度差异；层数不是本质，scale 是否沿
$K$ 变化才是。}
\label{fig:40b-format}
\end{figure}

\begin{figure}[H]
\centering
\begin{tikzpicture}[font=\scriptsize,
  stage/.style={draw, rounded corners=2pt, align=center, inner sep=3pt,
                minimum height=1.5cm, text width=3.2cm},
  arr/.style={->, thick}]
  \node[stage, fill=tkGreen!18] (main) at (1.85,1.4)
    {\textbf{(1) E4M3 mainloop}\\ $tCrC \mathrel{+}= \bar q_A\bar q_B^{\top}$\\
     两 stage TMA 流水\\（无任何 scale）};
  \node[stage, fill=tkOrange!25] (scale) at (5.55,1.4)
    {\textbf{(2) apply\_scales}\\ $tCrC \leftarrow \mathrm{fmaf}(tCrC\cdot\mathrm{sa}[m],$\\
     $\mathrm{sw}[n],\mathrm{bias}[n])$\\ 先乘 scale 的次序在此固化};
  \node[stage, fill=tkBlue!18] (lora) at (9.25,1.4)
    {\textbf{(3) gemm\_bf16\_lora}\\ $tCrC \mathrel{+}= D\,L_1^{\top}$，rank 切片 32 宽\\
     tid$<$32 预取前 2 片；rank 128 按需 refill};
  \node[stage, fill=tkRed!18] (out) at (12.95,1.4)
    {\textbf{(4) 单次舍入与写回}\\ NumericArrayConverter 一次\\
     FP32$\to$BF16 $\to$ STSM $\to$ TMA store};
  \draw[arr] (main) -- (scale);
  \draw[arr] (scale) -- (lora);
  \draw[arr] (lora) -- (out);
  \node[anchor=north west, align=left, text width=15.2cm, font=\tiny] at (0,-0.15)
    {四段共用同一个 \texttt{tCrC}（FP32 fragment）：(1)(3) 的两个 MMA 的 C partition
     被 static\_assert（gemm\_w8a8.cuh L97--102）锁成逐元素一致，跨乘积累加不需要
     retile；(2) 若挪到 (3) 之后，按式~\eqref{eq:40b-leak}，scale 会整体乘进低秩项
     $G$。(4) 是全 kernel 唯一一次 FP32 到 BF16 转换（L388--405）。};
\end{tikzpicture}
\caption{FP32 accumulator 的四段生命周期。scale 施加夹在两个乘积之间，
是这个 kernel 与普通 FP8 GEMM 唯一不能调换次序的地方。}
\label{fig:40b-acc-fusion}
\end{figure}

\section{性能分析}\label{sec:40b-perf}

本节数字全部取自 cache-dit PR \#1384 的验收记录（RTX PRO 5000，SM120，
2026-11-20；CUDA Event median），\textbf{本书未复测}。复测入口为
\texttt{bench/kernels/bench\_svdquant\_fp8.py}（默认 warmup 10、iters 30）。
\textbf{比值口径提醒}：下文所有"相对纯 FP8"的比值 $>1$ 表示\textbf{更慢}
（与第~\ref{ch:40}~章"相对 native 更快"的表相反，读表先看表头）。

\subsection{算子级：rank 越高越慢，且慢得有规律}

\begin{table}[H]
\centering\small
\caption{SVDQ-FP8 算子级耗时（$M{=}4096,N{=}12288,K{=}3072$，median CUDA
Event，含 producer 与 GEMM；数据取自 PR \#1384，本书未复测）。最后一列为
相对纯 FP8 的比值，$>1$ 表示更慢。}
\label{tab:40b-op}
\begin{tabular}{@{}lrrrr@{}}
\toprule
配置 & producer (ms) & GEMM (ms) & total (ms) & 相对纯 FP8 \\
\midrule
纯 FP8 per-row & --- & --- & 0.746 & $1.000\times$ \\
SVDQ-FP8 r32  & 0.046 & 0.781 & 0.827 & $1.108\times$ \\
SVDQ-FP8 r64  & 0.057 & 0.817 & 0.874 & $1.171\times$ \\
SVDQ-FP8 r128 & 0.074 & 0.851 & 0.925 & $1.239\times$ \\
\bottomrule
\end{tabular}
\end{table}

三个读表点。其一，开销大头在 GEMM 侧不在 producer：producer 只占
$0.046$--$0.074\,\mathrm{ms}$（含 pybind 层为 split-K $D$ 分配的
\texttt{at::zeros} scratch），随 rank 增长缓慢（amax 与量化写回与 rank
无关，随 rank 增长的只有 LoRA-down 分支：\texttt{lora\_down} 的装载与
BF16 MMA）。其二，GEMM 侧的增量是 LoRA-up：$\mathrm{nRt}=\mathrm{rank}/32$
在 rank 32/64/128 上取 1/2/4，BF16 MMA 累加与 cp.async 填环的字节量随
rank 线性增长，$0.781\to0.851\,\mathrm{ms}$ 的 $9\%$ 斜率与之吻合。
其三，"慢"是设计内代价——这张表买的是
表~\ref{tab:40b-acc} 的精度列。

\subsection{端到端：FLUX.1-dev}

\begin{table}[H]
\centering\small
\caption{FLUX.1-dev 端到端（默认 prompt/seed 与生成设置；SVDQ 预处理用
\texttt{--svdq-quantize-device cuda}，只计推理时间；数据取自 PR \#1384，
本书未复测）。最后一列为相对纯 FP8 的比值，$>1$ 表示更慢。}
\label{tab:40b-e2e}
\begin{tabular}{@{}llrrr@{}}
\toprule
模式 & 配置 & 推理 (s) & 相对 BF16 & 相对纯 FP8 \\
\midrule
eager   & BF16         & 17.00 & $1.00\times$ & --- \\
eager   & FP8 per-row  &  9.49 & $1.79\times$ & $1.00\times$ \\
eager   & SVDQ-FP8 r32 & 10.07 & $1.69\times$ & $1.06\times$ \\
eager   & SVDQ-FP8 r64 & 10.26 & $1.66\times$ & $1.08\times$ \\
eager   & SVDQ-FP8 r128& 10.62 & $1.60\times$ & $1.12\times$ \\
compile & BF16         & 14.85 & $1.00\times$ & --- \\
compile & FP8 per-row  &  7.14 & $2.08\times$ & $1.00\times$ \\
compile & SVDQ-FP8 r32 &  8.08 & $1.84\times$ & $1.13\times$ \\
compile & SVDQ-FP8 r64 &  8.29 & $1.79\times$ & $1.16\times$ \\
compile & SVDQ-FP8 r128&  8.65 & $1.72\times$ & $1.21\times$ \\
\bottomrule
\end{tabular}
\end{table}

e2e 把算子级的百分比稀释了约一半（Linear 之外的时间摊薄）：算子级慢
$11\%$--$24\%$，到 e2e 只剩 $6\%$--$21\%$；相对 BF16 仍保持
$1.60$--$1.84\times$ 的加速。也就是说，W8A8 落在"比 BF16 快近
一倍、比纯 FP8 慢一成"的窄带里，它的价值完全由精度表决定。

\subsection{精度：rank 128 才拉开差距}

\begin{table}[H]
\centering\small
\caption{FLUX.1-dev 精度（PSNR/SSIM 对各自 BF16 参考图；单 prompt/seed
参考测量，非统计质量门；数据取自 PR \#1384，本书未复测）。}
\label{tab:40b-acc}
\begin{tabular}{@{}llrr@{}}
\toprule
模式 & 配置 & PSNR (dB) & SSIM \\
\midrule
eager   & FP8 per-row  & 23.06 & 0.9215 \\
eager   & SVDQ-FP8 r32 & 22.18 & 0.9175 \\
eager   & SVDQ-FP8 r64 & 22.33 & 0.9202 \\
eager   & SVDQ-FP8 r128& \textbf{28.66} & \textbf{0.9713} \\
compile & FP8 per-row  & 22.07 & 0.9218 \\
compile & SVDQ-FP8 r32 & 21.63 & 0.9172 \\
compile & SVDQ-FP8 r64 & 21.05 & 0.9057 \\
compile & SVDQ-FP8 r128& \textbf{27.23} & \textbf{0.9645} \\
\bottomrule
\end{tabular}
\end{table}

结论分三层，每层都带上 PR 的限定语。第一，\textbf{rank 128 是唯一的
精度档}：eager $+5.60\,\mathrm{dB}$、compile $+5.16\,\mathrm{dB}$，SSIM
同步抬升约 $0.04$--$0.05$，PR 的目视检查（无结构损坏、无幻影物体、无
文字丢失）也指向 r128 最接近 BF16 参考。这个量级的差距远大于表内其它
任何两组之间的差，方向性可信。第二，\textbf{rank 32/64 不是精度档}：
四组比较全部略低于纯 FP8（$-0.44$--$-1.02\,\mathrm{dB}$），与作者
"The smaller ranks are not claimed to improve every end-to-end sample over
plain FP8" 一致；这两档只能理解为"为更高 rank 探路"或显存受限的妥协。
第三，\textbf{小于 $1\,\mathrm{dB}$ 的差距不下结论}：PR 声明 randomized
SVD 与 split-K 原子累加会造成 run-to-run 轨迹漂移，本书也未建立 SVDQ
的 PSNR 噪声带，r32/r64 与纯 FP8 之间的零点几分贝既可能是实现噪声、
也可能是真实小劣化。图~\ref{fig:40b-tradeoff} 把"花多少、买多少"画在
一张平面上：横轴向右是 rank 升高带来的耗时增加，纵轴向上是 PSNR，只有
r128 的点真正抬出了精度的台阶。

\begin{figure}[H]
\centering
\begin{tikzpicture}[font=\scriptsize,
  axis/.style={->, thick},
  bar/.style={draw, fill=tkBlue!25},
  dot/.style={circle, draw=tkRed!75!black, fill=tkRed!40, inner sep=1.6pt}]
  % 坐标轴：横轴 x=(t-9.4)*8.15（9.4..10.7 s -> 0..10.6），纵轴 y=(PSNR-21.5)/2
  \draw[axis] (0,0) -- (10.6,0);
  \node[below, font=\tiny] at (5.3,-1.05) {eager 推理耗时 (s) $\to$ 更慢};
  \draw[axis] (0,0) -- (0,4.4);
  \node[rotate=90, font=\tiny] at (-0.85,2.2) {PSNR (dB) $\to$ 更高};
  \foreach \x/\t in {0.73/9.49, 5.46/10.07, 7.01/10.26, 9.94/10.62}
    \node[below, font=\tiny] at (\x,0) {\t};
  \foreach \y/\t in {0.25/22, 1.25/24, 2.25/26, 3.25/28}
    \node[left, font=\tiny] at (0,\y) {\t};
  % 柱高 = (PSNR-21.5)/2：23.06 22.18 22.33 28.66 -> 0.78 0.34 0.415 3.58
  \draw[bar] (0.38,0) rectangle (1.08,0.78);
  \draw[bar] (5.11,0) rectangle (5.81,0.34);
  \draw[bar] (6.66,0) rectangle (7.36,0.415);
  \draw[bar] (9.59,0) rectangle (10.29,3.58);
  % 点（PSNR 值）
  \node[dot] (p0) at (0.73,0.78) {};
  \node[dot] (p1) at (5.46,0.34) {};
  \node[dot] (p2) at (7.01,0.415) {};
  \node[dot] (p3) at (9.94,3.58) {};
  \node[right=2pt, font=\tiny, text=tkRed!75!black] at (p0) {23.06};
  \node[right=2pt, font=\tiny, text=tkRed!75!black] at (p1) {22.18};
  \node[right=2pt, font=\tiny, text=tkRed!75!black] at (p2) {22.33};
  \node[left=2pt, font=\tiny, text=tkRed!75!black] at (p3) {28.66};
  % 配置标签
  \node[below, font=\tiny, align=center] at (0.73,-0.62) {纯 FP8\\per-row};
  \node[below, font=\tiny] at (5.46,-0.50) {r32};
  \node[below, font=\tiny] at (7.01,-0.50) {r64};
  \node[below, font=\tiny] at (9.94,-0.50) {r128};
\end{tikzpicture}
\caption{W8A8 的精度--耗时平面（eager 档；数据取自 PR \#1384，本书未复测）。
柱高与红点都是 PSNR（左轴刻度），横轴按表~\ref{tab:40b-e2e} 的耗时线性
摆放；纵轴自 $21.5\,\mathrm{dB}$ 起，柱高是截断后的差值。r32/r64 的点贴在
纯 FP8 之下，只有 r128 用 $+1.13\,\mathrm{s}$（相对纯 FP8）的耗时换来了
$+5.6\,\mathrm{dB}$ 的台阶。}
\label{fig:40b-tradeoff}
\end{figure}

与第~\ref{ch:40}~章 W4A4 的口径关系：W4A4 的性能叙事是"量化把 Linear
变便宜"（权重视显存 $4.5$ bit/元素，GEMM 两端吃低比特码），W8A8 没有
这条叙事，它只有"相对纯 FP8 多花 $11\%$--$24\%$ 算子时间、相对 BF16
仍快 $1.6$--$1.8\times$"的窄带——两章的加速比口径不可互相推导，也不
可拼接。

\section{实践经验}

\begin{itemize}[itemsep=1pt,topsep=2pt]
  \item \textbf{复用面比新增面大}。分解、平滑、LoRA 双分支数据流、
        epilogue 单次舍入协议，全部沿用 W4A4 的既有件；W8A8 真正新增的
        只有 amax kernel、split-K 的 $D$ 归约、以及 GEMM 里的 LoRA
        预取/refill 调度。从 W4A4 扩展到 FP8 的工作量主要花在
        \emph{删掉} block-scaled 那套布局上，而不是加新东西。
  \item \textbf{per-token 尺度的两个精度细节}。其一是 satfinite：E4M3
        转换溢出时钉在最大有限值，不刷 inf，下游 FP32 累加不会被污染；
        其二是倒数预计算：CTA 开头把 128 行的 $1/\mathrm{sa}$ 放进 smem，
        量化循环里对 sa 的除法变成乘法（逐输入通道的平滑仍是一次除法）；
        源码注释（quantize\_w8a8.cuh L59--60）特意
        交代，倒数乘法与 fast-math 除法可能让 E4M3 翻一个 ULP，所以转换
        指令固定为 RTN 加 satfinite 的
        \texttt{\_\_nv\_cvt\_float2\_to\_fp8x2}，与原生 FP8 producer 逐
        指令一致；想跨实现对拍数值时，这一条比任何容差都先决。
  \item \textbf{两档几何约束并存的原因}：producer 的 K tile 是 64
        （\texttt{kW8A8QuantBK}），所以 $K$ 只需 64 的倍数；GEMM 的
        K tile 是 128，所以合成链要求 $K$ 是 128 的倍数。模块路径下
        非 128 整数倍的几何由 buffer pad 与 \texttt{F.pad} 静默吸收，直调
        算子才会吃到 \texttt{ValueError}——排错时先确认走的是哪条路径。
  \item \textbf{rank 0 是合法的退化档}：低秩张量传 None 即退化为纯
        per-token FP8，producer 总是按 split 目标调度（没有 $D$ 需要归约，
        split 只为多造 CTA 掩延迟）。这给"同一 kernel 家族覆盖纯 FP8 与
        SVDQ"留下了一个开关位。
  \item \textbf{TMA 描述符缓存的 key 是}指针 $+$ shape $+$ rank
        （上限 32 条，超出整体清空）。描述符只是拷贝配方、不含数据，
        缓存它没有任何数值风险；但这也意味着换张量不换指针的复用
        场景通常会命中，排障时不要把它当成副作用怀疑对象。
  \item \textbf{显式拒绝优于静默回退}：序列化、few\_shot、非默认
        svdq\_kwargs、非 BF16 dtype 全部硬报错。W8A8 是精度档，用它
        的人对"配置没有生效"最敏感：静默忽略一个 knob 比直接报错
        的代价高得多。
\end{itemize}

\section{常见坑}

\begin{itemize}[itemsep=1pt,topsep=2pt]
  \item \textbf{把 W8A8 当压缩方案}：权重 8 bit/元素，权重视显存约为
        NVFP4 pack 的 $1.8$ 倍；它没有体积收益，只有精度收益。
  \item \textbf{用 W4A4 的格式存 FP8 artifact}：W4A4 的 save/load 描述
        的是 NVFP4 pack 布局，与 E4M3 残差互不兼容，代码已用
        \texttt{NotImplementedError} 焊死；自己写序列化前先想清楚要存
        的是哪一套约定。
  \item \textbf{把 scale 后乘到 dense 输出}：式~\eqref{eq:40b-leak} 的
        泄漏项是 $(\mathrm{sa}\cdot\mathrm{sw}-1)\,G$，量纲上 sa、sw 都
        是 absmax/448 量级，低秩项会被整体放大；\texttt{apply\_scales}
        必须先于 \texttt{gemm\_bf16\_lora}，这是次序约束不是风格偏好。
  \item \textbf{混淆两条消费链}：$\bar q_A$ 吃 $X/s$、$D$ 吃原始 $X$
        （lora\_down 自带 $1/s$）。两边都"碰一次 s"但位置不同，
        想重排任何一边前先按式~\eqref{eq:40b-offline} 核对量纲。
  \item \textbf{拿 cuh 模板默认值推 smem}：\texttt{kStages\_} 的默认是
        3，实际实例化是 \texttt{<128,2,128>}；96 KiB 的预算只按两
        stage 成立；改 stage 数前先重算这条静态断言。
  \item \textbf{读 r32/r64 的 PSNR}：PR 表里小 rank 略低于纯 FP8 是
        数据本身如此，不是实现坏了；$1\,\mathrm{dB}$ 以内的差距在
        run-to-run 漂移面前不下结论。
  \item \textbf{把 refill 的条件记成 rank $\ge$ 64}：触发条件是
        $\mathrm{rt}\ge\mathrm{kLoraStages}=2$，即 $\mathrm{nRt}>2$，
        \textbf{只有 rank 128}（nRt $=4$ 的下标 2、3，即第 3、4 片）会走到
        NamedBarrier 保护的按需补片路径。
  \item \textbf{E4M3 转换的 1 ULP 分歧}：普通转换路径与
        RTN$+$satfinite 路径在个别值上差一个码点；producer 为此固定
        用后者。任何"先除再转"的重写都可能引入跨实现的 1 ULP 级
        不一致。
\end{itemize}

\section{面试要点速查}

\begin{itemize}[itemsep=1pt,topsep=2pt]
  \item Q：SVDQuant 扩展到 FP8 需要什么数学条件？\\
        A：分解 $W s=L_1L_2+R$ 与误差上界对任意反量化映射 $\mathcal Q$
        成立（第~\ref{ch:38}~章 Proposition~4.1），把 $\mathcal Q$ 换成
        per-row E4M3 映射后命题链原样有效（\S\ref{sec:40b-motivation}~节）。
  \item Q：W8A8 与 W4A4 在 mainloop 上的本质差异是什么？\\
        A：不是码点位宽，而是 \textbf{scale 是否沿 $K$ 变化}：W4A4 的
        micro-scale 沿 $K$ 分布、提不出归约，必须 block-scaled MMA；
        W8A8 的行/列尺度沿 $K$ 不变，scale 全推到 epilogue 一次
        \texttt{fmaf}（式~\eqref{eq:40b-fold}，
        表~\ref{tab:40b-compare}）。
  \item Q：两个不同 MMA 怎么共享一个 FP32 accumulator？\\
        A：AtomLayout 与 M/N permutation 同构使 C partition 逐类型相等
        （gemm\_w8a8.cuh L97--102 的 static\_assert），两个
        \texttt{cute::gemm} 直接写同一个 \texttt{tCrC}，无需 retile
        （\S\ref{sec:40b-gemm}~节）。
  \item Q：为什么 scale 必须先乘再 LoRA 累加？\\
        A：后乘会把 sa$\cdot$sw 乘进低秩项，误差
        $(\mathrm{sa}\cdot\mathrm{sw}-1)\,G$（式~\eqref{eq:40b-leak}）；
        W4A4 靠 LoRA-first 的代数补偿规避，W8A8 的低秩项是原始量纲，
        只能由调用顺序保证。
  \item Q：producer 为什么拆成 amax 与主 kernel 两个 launch？\\
        A：per-token 尺度要等整行 amax，行级 grid（每行一 CTA）才有
        足够的独立 load 流掩 HBM 延迟；128 行一 tile 的 CTA 自己做
        不到（\S\ref{sec:40b-producer}~节）。
  \item Q：split-K 的 $D$ 为什么用 relaxed atomic 加 finalize？\\
        A：并行度换可复现性：多个 K 分片的 FP32 partial 原子归进
        scratch，finalize kernel 做唯一的 BF16 舍入；单次舍入保住了，
        逐位可复现放弃，等价性用误差门槛判。
  \item Q：rank 0 是干什么的？\\
        A：退化档：低秩张量传 None 即纯 per-token FP8，同一份 kernel 家族
        覆盖"无低秩修正"的基线。
\end{itemize}

\section{最小测试与动手实验}

下面的入口要在完整的 \texttt{fdf42708} 树上执行（附录~\ref{app:C} 的混合
checkout 只服务源码引用）。跑之前先重编扩展：

\begin{lstlisting}[style=console]
CACHE_DIT_BUILD_CUDA=1 CACHE_DIT_CUDA_ARCH_LIST=120a \
  bash tools/build_fast.sh install -e . --no-build-isolation --no-deps
\end{lstlisting}

\noindent PR \#1384 另报既有 SVDQ 回归 255 passed / 6 skipped。

\begin{lstlisting}[style=console]
# 读者侧验证（cache-dit 仓库，SM120 卡；本书未复跑）：
$ python -m pytest tests/kernels/test_svdquant_fp8.py \
                   tests/kernels/test_ops_dispatch.py \
                   tests/quantization/test_quantize_config.py -q
# PR #1384 自报这三个文件合计 164 passed；test_svdquant_fp8.py 单文件
# 含 18 个 test 函数（rank 0/32/64/128 x 多种 shape 参数化），覆盖数值
# 契约、split-K 路径、fake/compile 与拒绝用例。

# 性能复测（RTX PRO 5000，SM120）：
$ python bench/kernels/bench_svdquant_fp8.py        # 默认 warmup 10 / iters 30

# 动手实验：
# 1) rank 扫描：固定模型与 prompt，rank = 0/32/64/128 出精度-耗时曲线，
#    验证"只有 r128 抬出台阶"是否在你的模型上复现；
# 2) 关 LoRA：把 rank=0 的退化档与纯 FP8 对拍，分离"低秩分支"与
#    "per-token 量化本身"的贡献。
\end{lstlisting}

\section*{延伸阅读}

\begin{itemize}[itemsep=1pt,topsep=2pt]
  \item cache-dit PR \#1384（\texttt{vipshop/cache-dit}）：本章全部性能与
        精度数字的出处，含 eager/compile 两组对比图；
  \item 第~\ref{ch:38}~章：SVDQuant 数学主线（Proposition~4.1 与
        Eckart--Young 是 \S\ref{sec:40b-motivation} 的两处引用）；
  \item 第~\ref{ch:39}--\ref{ch:40}~章：W4A4 的 producer 与 GEMM，本章
        的对照基线；
  \item 第~\ref{ch:41}~章：原生 FP8（无低秩分支）的 CUTLASS Collective
        实现，式~\ref{eq:41-fp8-evt} 的 EVT 形态可与
        式~\eqref{eq:40b-fold} 对照；
  \item 第~\ref{ch:31}~章：NVFP4 两级 scale 的格式定义，理解
        表~\ref{tab:40b-compare} 的 scale 列的前提。
\end{itemize}
